\section{Methods}

%% ---- Data Sources & Acquisition --------------------------------
\subsection{Data Sources \& Acquisition}

\subsubsection{Data Source}

All data were collected from \textbf{CoinGecko} historical data pages, which provide daily OHLCV figures going back several years.
The target URL pattern is:
\begin{center}
  \texttt{https://www.coingecko.com/en/coins/\{coin\_id\}/historical\_data}
\end{center}

\subsubsection{Scraper Architecture}

The scraper is built on \textbf{Scrapy} with the \texttt{scrapy-playwright} extension, which drives a headless Chromium browser to handle JavaScript-rendered tables.
Key configuration parameters are listed in Table~\ref{tab:scraper-config}.

\begin{table}[H]
  \centering
  \caption{Scrapy / Playwright configuration.}
  \label{tab:scraper-config}
  \begin{tabular}{@{}L{4cm}L{3.5cm}L{6cm}@{}}
    \toprule
    \textbf{Setting} & \textbf{Value} & \textbf{Purpose} \\
    \midrule
    \texttt{DOWNLOAD\_DELAY}       & 2.5 s  & Polite crawl rate \\
    \texttt{AUTOTHROTTLE\_ENABLED} & True   & Adaptive rate limiting \\
    \texttt{PLAYWRIGHT\_BROWSER}   & chromium & Headless JS rendering \\
    \texttt{ROBOTSTXT\_OBEY}       & True   & Respects robots.txt \\
    \bottomrule
  \end{tabular}
\end{table}

\subsubsection{Parsing Logic}

For each table row the spider reads four columns: date, market cap, 24 h volume, and closing price.
Because CoinGecko does not expose an open price in the historical table, \texttt{price\_change\_pct} is always yielded as \texttt{None} and recomputed during preprocessing as a day-over-day return.

The spider applies a three-year date filter and repeatedly clicks the \emph{Show More} button (up to 80 times) to load the full history.
On close it logs the start date, end date, and item count per coin for audit purposes.

\subsubsection{Raw Data Schema}

Each raw CSV contains the fields shown in Table~\ref{tab:raw-schema}.

\begin{table}[H]
  \centering
  \caption{Raw data schema (\texttt{data/raw/\{coin\_id\}.csv}).}
  \label{tab:raw-schema}
  \begin{tabular}{@{}L{3.5cm}L{2.5cm}L{7cm}@{}}
    \toprule
    \textbf{Field} & \textbf{Type} & \textbf{Description} \\
    \midrule
    \texttt{date}             & string        & ISO 8601 (\texttt{YYYY-MM-DD}) \\
    \texttt{coin\_id}          & string        & CoinGecko slug \\
    \texttt{coin\_name}        & string        & Human-readable name \\
    \texttt{symbol}           & string        & Ticker (e.g.\ BTC) \\
    \texttt{price\_usd}        & float         & Daily closing price in USD \\
    \texttt{market\_cap\_usd}  & float         & Market capitalisation in USD \\
    \texttt{volume\_24h\_usd}  & float         & 24-hour trading volume in USD \\
    \texttt{price\_change\_pct}& float         & Day-over-day \% change (recomputed) \\
    \texttt{scraped\_at}       & string (UTC)  & Timestamp of scrape \\
    \bottomrule
  \end{tabular}
\end{table}

%% ---- Data Preprocessing ----------------------------------------
\subsection{Data Preprocessing}

\subsubsection{Overview}

Raw CSVs are read from \texttt{data/raw/} and passed through a five-step cleaning pipeline implemented in \texttt{src/preprocessing.py}.
The cleaned output is written to \texttt{data/processed/\{coin\_id\}\_cleaned.csv}.

\subsubsection{Processing Steps}

\textbf{Step 1 — Type casting.}
The \texttt{date} column is parsed to \texttt{datetime64[ns]}; the four numeric fields are cast to \texttt{float64}.

\textbf{Step 2 — Deduplication.}
Rows are sorted by \texttt{scraped\_at} descending; duplicates on the composite key \texttt{(coin\_id, date)} are dropped, retaining the most recent scrape.

\textbf{Step 3 — Missing value handling.}
The strategy varies by field importance (Table~\ref{tab:missing}).

\begin{table}[H]
  \centering
  \caption{Missing value handling strategy.}
  \label{tab:missing}
  \begin{tabular}{@{}L{3.5cm}L{9cm}@{}}
    \toprule
    \textbf{Field} & \textbf{Strategy} \\
    \midrule
    \texttt{price\_usd}        & Drop row — price is the core signal and cannot be fabricated \\
    \texttt{market\_cap\_usd}  & Forward-fill within coin group; flag with \texttt{market\_cap\_usd\_imputed = True} \\
    \texttt{volume\_24h\_usd}  & Forward-fill within coin group; flag with \texttt{volume\_24h\_usd\_imputed = True} \\
    \texttt{price\_change\_pct}& Recompute as $p_t / p_{t-1} - 1$; prefer scraped value when present \\
    \bottomrule
  \end{tabular}
\end{table}

\textbf{Step 4 — Outlier flagging.}
Rows where $|\texttt{price\_change\_pct}| > 50\%$ are marked \texttt{is\_outlier = True}.
They are \emph{retained} rather than dropped: extreme price swings such as exchange listings or market crashes are real events that should not be silently removed.

\textbf{Step 5 — Sorting.}
The output is sorted by \texttt{(coin\_id, date)} ascending.

%% ---- Feature Engineering ---------------------------------------
\subsection{Feature Engineering}

\subsubsection{Overview}

Feature engineering is implemented in \texttt{src/features.py}.
All features are computed \emph{per coin} (grouped by \texttt{coin\_id}) to prevent cross-coin leakage, then \emph{lagged by one day} so the model only sees information available before the prediction date.

\subsubsection{Computed Features}

\begin{table}[H]
  \centering
  \caption{Engineered features and lag status.}
  \label{tab:features}
  \begin{tabular}{@{}L{3cm}L{6.5cm}C{2cm}@{}}
    \toprule
    \textbf{Feature} & \textbf{Formula} & \textbf{Lagged?} \\
    \midrule
    \texttt{daily\_return}   & $p_t / p_{t-1} - 1$                            & Yes \\
    \texttt{ma\_7}           & 7-day rolling mean of \texttt{price\_usd}       & Yes \\
    \texttt{ma\_30}          & 30-day rolling mean of \texttt{price\_usd}      & Yes \\
    \texttt{volatility\_7}   & 7-day rolling std of \texttt{daily\_return}     & Yes \\
    \texttt{vol\_change}     & $v_t / v_{t-1} - 1$                             & Yes \\
    \texttt{price\_direction}& $1$ if \texttt{daily\_return} $> 0$, else $0$   & \textbf{No} (target) \\
    \bottomrule
  \end{tabular}
\end{table}

\subsubsection{Lag Semantics}

The prediction task is illustrated schematically below:

\[
  \underbrace{\texttt{daily\_return}_{t-1},\ \texttt{ma\_7}_{t-1},\ \texttt{ma\_30}_{t-1},\ \texttt{volatility\_7}_{t-1},\ \texttt{vol\_change}_{t-1}}_{\text{features from day }t-1}
  \;\longrightarrow\;
  \underbrace{\texttt{price\_direction}_{t}}_{\text{target at day }t}
\]

The first row per coin always has \texttt{NaN} for all lagged features (nothing precedes day 1) and is dropped before modeling.

%% ---- Models ----------------------------------------------------
\subsection{Models}

\subsubsection{Task \& Split Strategy}

The task is binary classification: predict \texttt{price\_direction} $\in \{0, 1\}$ on day $t$.
A \textbf{chronological train/test split} (80\,\% train, 20\,\% test) is applied \emph{per coin independently}; no shuffling is performed.
Shuffling would allow the model to see future data during training, producing optimistic metrics that do not reflect real-world performance.

\subsubsection{Model Descriptions}

Three models are evaluated (Table~\ref{tab:models}):

\begin{table}[H]
  \centering
  \caption{Models evaluated.}
  \label{tab:models}
  \begin{tabular}{@{}L{3.5cm}L{4.5cm}L{5cm}@{}}
    \toprule
    \textbf{Key} & \textbf{Class} & \textbf{Notes} \\
    \midrule
    Logistic Regression & \texttt{LogisticRegression(max\_iter=1000)} & Linear baseline \\
    Decision Tree       & \texttt{DecisionTreeClassifier()}           & Default depth, no pruning \\
    Linear Regression   & \texttt{LinearRegression()}                 & Continuous output thresholded at 0.5 \\
    \bottomrule
  \end{tabular}
\end{table}

\subsubsection{Evaluation Metrics}

Each model is scored on four metrics: accuracy, precision, recall, and F1-score (harmonic mean of precision and recall).
Accuracy is the primary metric for this study; the secondary metrics reveal class-level performance and potential imbalance effects.
